\documentclass[11pt]{article}
\usepackage{fontspec,amsmath,amssymb,unicode-math,booktabs,longtable,array,fvextra,hyperref,bookmark}
\usepackage{komorebi-publication}
\renewcommand{\KomorebiTitle}{Universal Residual Seeds}
\renewcommand{\KomorebiSubtitle}{Deformation forests, transport categories, and conditional system completion}
\renewcommand{\KomorebiSeries}{Core series X}
\renewcommand{\KomorebiKind}{Lean and evidence guide}
\renewcommand{\KomorebiShortTitle}{Paper X}
\renewcommand{\KomorebiTitleSize}{36}
\renewcommand{\KomorebiTitleLeading}{41}
\renewcommand{\KomorebiEdition}{Integrated conditional reading edition / 8 October 2026}
\KomorebiSetup

\newcommand{\CodeSlash}{\ensuremath{\big/}}
\begin{document}
\KomorebiCover
\begin{abstract}
This guide records the final conditional occupied-system completion endpoint for Paper X, its independent target-family and carrier-recovery outputs, its explicit semantic inputs, and its saved verification. The entire production module is printed from its source bytes. The earlier 20/20 conditional ledger and the later 19/20 unpromoted residual-system attempt are kept separate.
\end{abstract}
\tableofcontents
\section{The exact source contract}
The terminal module imports \nolinkurl{PaperX.HorizontalTransportCompletion}, \nolinkurl{PaperX.Claim91ConstructedClosure}, and \nolinkurl{PaperX.StrictCompatibleCarrierCompletion}. Its namespace is \nolinkurl{PaperX.ConditionalOccupiedSystemCompletion}.
\begin{longtable}{@{}p{.41\linewidth}p{.55\linewidth}@{}}\toprule
Declaration & Role\\\midrule
\nolinkurl{CarrierRecoveryAt} & Packages equivariant splitting, unique finite-trace recovery, continuity, Frobenius polynomial products and supplied local finite-copy comparisons.\\
\nolinkurl{finite_trace_separator_at} & Constructs good-prime separation from the arithmetic occurrence theorem at a distinguished coefficient place.\\
\nolinkurl{carrier_recovery_from_arithmetic_separator} & Passes the constructed separator to the compatible-carrier closure.\\
\nolinkurl{conditional_occupied_system_completion} & Conjoins transported target-family existence with pointwise carrier recovery.\\\bottomrule
\end{longtable}
\section{Read the final proof before its signature}
The final proof uses \nolinkurl{constructor}. Its first branch is the existing inverse-limit-closed horizontal completion theorem. Its second branch is \nolinkurl{carrier_recovery_from_arithmetic_separator}. This small assembly sits on substantial predecessor results; it does not reduce the semantic premises of those results.

The inputs include actual completed source members at each coefficient place, all-depth admissible class transport, an inverse-limit-closed target profile, and a finite pairwise-nonisomorphic compatible irreducible constituent family. Carrier and initial-subrepresentation presentations, multiplicities $m_i\leq c_i$, constituent Weil--Deligne and smooth-inertia data, and semilinear period comparisons are supplied. Additional local conditions must satisfy their declared finite-copy closure contract.

The target-family and carrier statements are independent conjuncts. The carrier result is at an arbitrary distinguished place. Applying it at every place requires the corresponding presentations there. Zero multiplicity is accepted. A nonzero occupied system requires separate nonzero inputs. The module-support alternative is inherited from predecessor source.
\section{The recorded verification}
The final release's \nolinkurl{VERIFICATION.json} records two successful warnings-as-errors builds from separate frozen-cache copies, equal corresponding proof objects, source scanning and a transitive axiom audit. Allowed axioms are \nolinkurl{propext}, \nolinkurl{Classical.choice}, and \nolinkurl{Quot.sound}. Controls reject false and placeholder proofs and omitted closure or separator premises. An injected custom axiom is detected after compilation.

The pinned toolchain is Lean 4.34.0-rc2, commit \nolinkurl{6a10ac8c22beadecabdbb0919c2b50214762f91d}; the Mathlib revision is \nolinkurl{aeb72f1e9981b6d35b426b16dca7759be3e7ba06}. The original verifier retains execution-workspace paths that must be adjusted when restoring the release. The compact preservation map and restore script retain every original leaf payload, including nested history archives as restorable directory content. They preserve content, not original ZIP container byte streams.

The October 5 audit reports compilation of 30 Paper X project modules including this endpoint. This publication checks its exact source hash and rebuilds the documents; it does not perform a fresh Lean build. The later stronger attempt's \nolinkurl{NO_PINNED_LEAN} is not a compiler rejection of this checked source.
\section{Production source, unchanged}
Source SHA-256: \nolinkurl{ba15fd1ea175fecfcc705243c75d1061c8bdf500221dafe86f29fe1c5d21c5e0}.
The following listing is the complete production module. Its UTF-8 file remains the executable source of record. Font substitution in this printed view is typographic only.
\VerbatimInput[fontsize=\scriptsize,breaklines=true,breakanywhere=true,numbers=left,numbersep=6pt,frame=lines,commandchars=\\\{\}]{ConditionalOccupiedSystemCompletion.render.txt}
\section{Companion records}
Read \nolinkurl{REPORT.md}, \nolinkurl{PROOF_AUDIT.md}, \nolinkurl{CLAIM_MAP.md}, \nolinkurl{CLAIM_STATUS.json}, \nolinkurl{REPRODUCE.md}, and the evidence files under \nolinkurl{formal_preservation/readable/checked_conditional/}. The later residual attempt and terminal reports are retained separately under \nolinkurl{formal_preservation/readable/later_residual_attempt/} and \nolinkurl{status_records/}.
\end{document}
